% The abstract: compared with the statements it summarizes before any freeze (the rule that
% summaries drop hypotheses).
\begin{abstract}
An image pyramid smooths and subsamples an image by a fixed ratio $q$ at every level. The classical
design literature constrains the kernel of one step and follows the equivalent kernels that the
pyramid builds level by level. Scale-space theory has a kernel at every scale and leaves the
subsampling outside its treatment. This paper studies kernel
families on $\RR^d$ under the hemigroup axioms of linear scale space, in which two consecutive
smoothing steps compose to one, with covariance under the powers of $q$ in place of every dilation.
It gives a description of the families by their \Levy{} data, a selection theorem and an
obstruction on the grid. The isotropic
families form three nested strata: all of them, those covariant under the powers of $q$, and those
covariant under every dilation. In the middle stratum the kernel from the scale $0$ to each scale
is semi-selfdecomposable: it is its own contraction by $q$, smoothed by an infinitely divisible law.
The scales are labelled so that
dilating the kernel at a scale $t$ by $q$ gives the kernel at $qt$, and an octave step goes from $t$ to $qt$. The ladder of octave steps between the stored levels determines the kernels at
the knots $q^n$. No law of bounded support other than a point mass is a step of any family that obeys
the cascade and is continuous in scale. Every octave step of the \Matern{} pyramid is one implicit
heat step, a Helmholtz resolvent, and the pyramid is covariant under every dilation. Up to a
relabelling $g$ of the scales with $g(qt) = q\,g(t)$, it is
the one isotropic family covariant under $q$ whose octave steps from every scale obey the maximum
principle of an implicit Euler step for every smooth rapidly decreasing input. Among symmetric
rotation-invariant laws the principle holds exactly for the Helmholtz resolvents, so Gaussian steps
fail it. The principle is required of the octave steps only. A step of the \Matern{} pyramid has it
only across a factor $q$ in scale, and no symmetric family has it at every step. A grid pyramid
applies one mask, a probability law, at every level, and it is called exact when the steps it
applies from level to level are the octave steps of one family continuous in scale and covariant
under $q$. On a lattice, a symmetric probability mask other than a point mass has an exact grid
pyramid if and only if it is infinitely divisible with a logarithmic moment. For laws the
sufficiency is a theorem of Maejima and Ueda, and continuity in scale gives the necessity. For an
integer ratio the kernels of an exact grid pyramid on $h\ZZ^d$ are singular. In
$d \ge 2$ no kernel at a positive scale of an exact grid pyramid whose mask is carried by a lattice
is rotation invariant. The same holds when the dilation is an invertible matrix, which may also turn
the lattice. The obstruction needs continuity in scale: some band-limited mask $w$ of infinite
support has a rotation-invariant law $\mu$ with $D_q\mu = \mu * w$. The lemmas, propositions and
theorems of Sections 3 to 6 and of the supplement have formal counterparts in Lean~4. All are machine-checked but one step and three
results proved in print only, in whole or in part: the proposition on matrix dilations, the
singularity clause on the kernels, and the band-limited example. The formal counterparts rest on
nine cited facts, taken as fourteen axioms that the supplement prints in Lean, and on two earlier
formal developments of the author's. One of the nine is a proposition of an earlier paper of the
author's that restates a theorem of Karlin. Neither the connecting prose nor the
numerical examples are machine-checked, and the formal statements and definitions were compared
with the printed ones by AI agents only.
\end{abstract}
